\documentclass[11pt,a4paper]{article}
\usepackage[margin=20mm,headheight=16pt]{geometry}
\usepackage{fontspec,kotex,amsmath,booktabs,longtable,array,xcolor,hyperref,fancyhdr,graphicx}
\setmainfont{DejaVu Serif}
\setsansfont{DejaVu Sans}
\setmonofont{DejaVu Sans Mono}
\setmainhangulfont{Noto Sans CJK KR}
\setmonohangulfont{Noto Sans CJK KR}
\definecolor{navy}{HTML}{173E59}
\definecolor{codebg}{HTML}{F3F6F8}
\hypersetup{colorlinks=true,linkcolor=navy,urlcolor=navy,pdftitle={GuardSynth Discussion}}
\urlstyle{same}
\setlength{\parindent}{0pt}
\setlength{\parskip}{4pt}
\setlength{\emergencystretch}{4em}
\setlength{\tabcolsep}{4pt}
\renewcommand{\arraystretch}{1.2}
\pagestyle{fancy}\fancyhf{}
\fancyhead[L]{\small GuardSynth}
\fancyhead[R]{\small Discussion}
\fancyfoot[C]{\thepage}
\renewcommand{\headrulewidth}{0.3pt}
\clubpenalty=10000\widowpenalty=10000
\raggedbottom
\begin{document}

\section*{Discussion}

\subsection*{A. The Role of Explicit Behavioral Constraints in CoC}

Whereas CoC connects a situation with reasons for an action, explicit behavioral constraints specify conditions that permit or restrict choices. The higher accuracy and lower violation rates of P1 and P2 relative to P0 across static, temporal, and maneuver tasks support this complementary role. Contract-pair evaluation is particularly informative: the same scene and candidates can require different choices under different rules. Constraint enhancement therefore supplements the explanation in CoC with explicit conditions for action selection, rather than replacing that explanation.\par

The weak performance of shuffled associations indicates that adding constraint sentences alone is insufficient; their correspondence with each case matters. This supports the methodological motivation to connect targets, regions, and times to observations while distinguishing observations from applicable rules. However, P0 does not receive the variable contract, so the observed gains include the effect of supplying additional information. They should not be generalized to an overall deficiency in CoC reasoning or to greater learning efficiency under identical information.\par

\subsection*{B. Formal Support for Constraint Development}

P1 and P2 are not competing external methods but different development paths within our constraint-enhancement approach. P1 directly supplies existing natural-language constraints; P2 specifies and verifies the corresponding rules in EBLC before generating CNL. The central value of formalization is not necessarily higher model performance. It is the ability to distinguish targets, scope, activation, maintenance and release conditions, restricted actions, and assumptions in a checkable representation. Natural-language constraints can also be reviewed by people; the EBLC path additionally makes specification, translation, and generated-language correspondence explicit objects of supported checks.\par

Behavioral evaluation and error injection support different aspects of this value. The former shows that CNL produced through the formalization path can support high action-selection performance; the latter demonstrates checks for omissions, boundary errors, and semantic mismatches. Differences between P1 and P2 varied across tasks and expressions, so universal equivalence or noninferiority is not established. The supported claim is that systematic constraint development and checking can accompany the benefits of constraint enhancement, not that verification itself causes the behavioral improvement.\par

\newpage
\subsection*{C. Evidence, Verification, and Compliance}

Reliability in this process has three distinct layers. Evidence validity concerns whether the selected observations, rules, and assumptions are appropriate for the scene. Representation and translation validity concern whether the specification and CNL express the intended content. Behavioral compliance concerns whether the model selects an action consistent with the supplied constraint. SMT verification checks logical properties under the specified model and assumptions, while supported correspondence checks examine specification-to-language agreement. Neither substitutes for establishing the factual validity of scene interpretation or rule evidence.\par

Undetected shared errors in evidence and specification expose a limitation of the first layer; residual model violations despite a correct specification expose a limitation of the third. These require different responses. Expected UNSAT when a prohibited action is conjoined with the contract supports its exclusion, whereas unexpected UNSAT for an admissible action motivates re-examining evidence, clauses, and translation. Specification verification and behavioral evaluation are thus complementary sources of evidence addressing different errors, rather than interchangeable assurances.\par

\subsection*{D. Scope and Further Development}

The behavioral evidence is limited to supplied synthetic rules, finite candidates, one VLM, and three optimization seeds. Candidate values are also provided in text, and some images are reused; performance on unseen values therefore does not establish visual generalization to new real-road scenes. P1 and P2 also receive constraints at evaluation, leaving internalization without supplied constraints unexamined. Sensitivity to paraphrasing describes the applicability of the model and its input interface; this study does not isolate underlying language capability as its sole cause. Error detection likewise remains bounded by supported syntax and injected error types.\par

Further work can evaluate how constraint applicability is established in real scenes and how generated CNL is used across models and environments. A promising division of work concentrates human review on observations, rules, and assumptions, while assigning structuring, specification checks, and language generation to machine processing. GuardSynth Studio can provide an implementation setting for this workflow, but extraction accuracy and reductions in authoring time require separate evaluation. The present study establishes a more limited foundation: evidence for constraint-enhanced CoC action selection and a systematic path from structured, checked constraints to natural-language learning inputs.\par
\end{document}